\scanpage{3}
% archive foliation: 2
\uline{Cor} Conditions équivalentes

 a) $C=C_\infty$ i.e. pour $\forall x\in C$, l'ens des longueurs des parties partant de $x$ est fini.

 b) $\exists\ C\xrightarrow{l}\mathbb N$ tel que si $y\in R(x)$, on ait $l(y)<l(x)$

[Alors $\forall x\in C$, si $l=l(x)$, on a $x\in C_l$]

\deleted{\unsure{Impl.}}

\uline{NB} Ces conditions impliquent que $\nexists$ partie infinie.
\deleted{Mais} Si $C$ est \deleted{\unsure{in}}fini, \deleted{\unsure{l'inverse est vrai}}\inserted{la réciproque est vraie} : en effet, tte partie \deleted{\unsure{est finie}}\inserted{a alors une longueur} \unsure{au plus} $(\operatorname{card}C)-1=N-1$, i.e. $C=C_{N-1}$

\deleted{\uline{Prop} Si $C$}

\textcircled{3} \uline{Stratégie pour $j\in J$}
\[
\Sigma\subset(C(j)\times C)\cap R\quad\text{tel que}
\]
\[
\forall x\in C(j),\quad\text{on a}\quad(R(x)\ne\varnothing\ (\text{i.e. }x\notin C_0)\Longrightarrow\Sigma(x)\ne\varnothing)
\]
Partie $\underline x=(x_0,\ldots,x_n)$ \uline{compatible} avec stratégie $\Sigma$ pour $j$ :
\[
\forall\,0\leq i\leq n-1\quad\text{tel que }x_i\in C(j),\quad\text{on a }x_{i+1}\in\Sigma(x_i)
\]
\margin{Idem pour \unsure{partie infinie|parties infinies}}

\uline{Stratégie gagnante} \inserted{(pour $j$)} pour $x$ $(\in C)$ :

 a) Toute partie partant de $x$, compatible avec $\Sigma$, est finie

(i.e. $\nexists$ partie infinie partant de $x$, comp. avec $\Sigma$)

 b) \inserted{pour} Toute partie terminée, partant de $x$, \inserted{$(x_0=x,x_1,\ldots,x_n)$} compatible avec $\Sigma$, on a $x_n\in G_0(j)$.

On dit que $x\in C$ est \uline{gagnante} \uline{pour} $j$ si $\exists$ stratégie \deleted{gagnante}\inserted{pour $j$} gagnante pour $x$.
L'ens de ces $x$ est noté $G(j)$

\uline{NB} $G(j)\cap C_0=G_0(j)$

Posons $\overline C(j)=C\setminus C(j)$, $\overline G(j)=C\setminus G(j)$\ldots
